\section{Training curves and split check}\label{app:training}
Figure~\ref{fig:training} shows development and held-out RMSE by epoch for the five seeds of the $d=64$ autoencoder; early stopping restored the best epoch (mean \bestEpoch). With a 70/10/20 split, in which early stopping uses a 10\% validation subset of the development players and the 20\% test players are never touched, the test RMSE is \rmseTestThreeWay{} (SD \rmseTestThreeWaySd{} over \nThreeWay{} seeds) against \rmseHeld{} in the main design.
\begin{figure}[h]\centering\includegraphics[width=0.55\textwidth]{figA1_training}
\caption{Development-set (red) and held-out (blue) reconstruction RMSE per epoch for five training seeds, $d=64$.}\label{fig:training}\end{figure}

\section{Single-run UMAP+HDBSCAN: selection and stability}\label{app:single}
\begin{table}[h]\centering\scriptsize
\caption{Configurations selected on the development set by three rules and their behaviour on held-out and refitted data. $^a$UMAP dimension / neighbours / minimum distance / HDBSCAN minimum cluster size / minimum samples. Sil.\ = Silhouette in the UMAP space, noise excluded.}
\label{tab:selection}
\resizebox{\textwidth}{!}{\begin{tabular}{llccccccccc}\toprule
 & & \multicolumn{4}{c}{Development (80\%)} & \multicolumn{3}{c}{Held-out (20\%)} & \multicolumn{2}{c}{Refit on all} \\ \cmidrule(lr){3-6}\cmidrule(lr){7-9}\cmidrule(lr){10-11}
Selection rule & Configuration$^a$ & $k$ & Noise \% & Sil. & DBCV & $k$ & Noise \% & Sil. & $k$ & Noise \% \\ \midrule
DBCV (prespecified) & 10/10/0.0/15/15 & 7 & 1 & 0.36 & 0.48 & 7 & 1 & 0.36 & 13 & 12 \\
Silhouette & 2/30/0.0/20/5 & 24 & 25 & 0.52 & 0.19 & 23 & 56 & 0.47 & 14 & 10 \\
Composite score (2025 version) & 2/15/0.0/40/5 & 14 & 13 & 0.43 & 0.17 & 14 & 29 & 0.42 & 7 & 3 \\
\bottomrule\end{tabular}}
\end{table}
\begin{table}[h]\centering\small
\caption{Reproducibility of single runs. ARI over UMAP seeds: mean (range) of pairwise ARI among ten UMAP random states, averaged over three autoencoder seeds. ARI over AE seeds: mean pairwise ARI among the three autoencoder seeds at a fixed UMAP seed. $^a$Configurations: \configList.}
\label{tab:stability}
\resizebox{\textwidth}{!}{\begin{tabular}{llccccc}\toprule
Config.$^a$ & Extraction & ARI, UMAP seeds (range) & ARI, AE seeds & $k$ range (median) & Noise \% & Largest cluster \% \\ \midrule
C1 & eom & 0.60 (0.03--0.95) & 0.38 & 2--19 (9) & 13 & 49 \\
C1 & leaf & 0.64 (0.50--0.75) & 0.23 & 15--20 (18) & 37 & 8 \\
C4 & eom & 0.66 (0.03--0.99) & 0.34 & 2--14 (8) & 5 & 63 \\
C4 & leaf & 0.32 (0.18--0.50) & 0.18 & 12--20 (16) & 28 & 9 \\
C2 & eom & 0.63 (0.03--0.95) & 0.31 & 3--18 (12) & 19 & 34 \\
C2 & leaf & 0.59 (0.44--0.71) & 0.25 & 14--19 (16) & 34 & 9 \\
C3 & eom & 0.44 (0.02--0.98) & 0.09 & 2--22 (18) & 20 & 33 \\
C3 & leaf & 0.57 (0.44--0.70) & 0.20 & 20--27 (23) & 35 & 6 \\
\bottomrule\end{tabular}}
\end{table}

\section{Consensus diagnostics, cross-representation agreement and macro-types}\label{app:k}
Table~\ref{tab:kcurve} lists PAC, consensus Silhouette and the share of stable players for every $k$. PAC decreases almost monotonically in $k$ for every representation, which is why it cannot select $k$ by itself and the consensus Silhouette was used; the PAC-minimum rule would select $k=12$ for every representation. Table~\ref{tab:crossrep} gives the agreement between representations at $k=8$, and Table~\ref{tab:profiles4} the rule-selected macro-types with their nesting of the fine types (Table~\ref{tab:nesting}).
\begin{table}[h]\centering\small
\caption{Consensus diagnostics by number of clusters $k$ for the summary-statistic representation and the five-seed autoencoder representation.}\label{tab:kcurve}
\begin{tabular}{ccccccc}\toprule
 & \multicolumn{3}{c}{Summary statistics} & \multicolumn{3}{c}{Autoencoder $d=64$, 5 seeds} \\ \cmidrule(lr){2-4}\cmidrule(lr){5-7}
$k$ & PAC & Sil.$_C$ & Stable \% & PAC & Sil.$_C$ & Stable \% \\ \midrule
2 & 0.04 & 0.98 & 100 & 0.41 & 0.80 & 97 \\
3 & 0.11 & 0.92 & 99 & 0.57 & 0.48 & 80 \\
4 & 0.08 & 0.93 & 99 & 0.57 & 0.39 & 69 \\
5 & 0.11 & 0.90 & 99 & 0.50 & 0.37 & 66 \\
6 & 0.17 & 0.83 & 98 & 0.42 & 0.38 & 65 \\
7 & 0.14 & 0.80 & 96 & 0.36 & 0.42 & 66 \\
8 & 0.07 & 0.90 & 99 & 0.32 & 0.44 & 67 \\
9 & 0.13 & 0.80 & 97 & 0.30 & 0.38 & 51 \\
10 & 0.17 & 0.61 & 88 & 0.27 & 0.36 & 49 \\
11 & 0.16 & 0.63 & 89 & 0.25 & 0.34 & 42 \\
12 & 0.15 & 0.59 & 87 & 0.23 & 0.34 & 42 \\
\bottomrule\end{tabular}
\end{table}
\begin{table}[h]\centering\small
\caption{Agreement (ARI) between the $k=8$ consensus partitions obtained from different representations of the same \nPlayers{} careers.}
\label{tab:crossrep}
\begin{tabular}{lccccccc}\toprule
 & Summary & Resampled & PCA & DTW & AE (5 seeds) & AE (seed 0) & AE+UMAP \\ \midrule
Summary & 1.00 &  &  &  &  &  &  \\
Resampled & 0.32 & 1.00 &  &  &  &  &  \\
PCA & 0.18 & 0.12 & 1.00 &  &  &  &  \\
DTW & 0.19 & 0.27 & 0.06 & 1.00 &  &  &  \\
AE (5 seeds) & 0.22 & 0.20 & 0.12 & 0.12 & 1.00 &  &  \\
AE (seed 0) & 0.20 & 0.18 & 0.12 & 0.13 & 0.57 & 1.00 &  \\
AE+UMAP & 0.19 & 0.15 & 0.11 & 0.04 & 0.29 & 0.20 & 1.00 \\
\bottomrule\end{tabular}
\end{table}
\label{app:nesting}
\begin{table}[h]\centering\scriptsize
\caption{Macro-types ($k=\mainK$, rule-selected) of the summary-statistic consensus. Stability = median player stability; PER, seasons and slope are cluster medians.}
\label{tab:profiles4}
\resizebox{\textwidth}{!}{\begin{tabular}{lcccccccclp{3.6cm}}\toprule
Type & $n$ & Stability & Mean PER & Peak PER & Seasons & Slope & All-Star \% & Career WS & Label (rule) & Nearest to medoid \\ \midrule
T1 & 369 & 0.98 & 16.5 & 20.4 & 12 & -0.30 & 52 & 52.5 & moderate-PER, long career, declining & Josh Smith, Amir Johnson, Eric Bledsoe \\
T2 & 560 & 0.98 & 12.5 & 16.2 & 11 & -0.34 & 8 & 24.7 & moderate-PER, long career, declining & Mike Sanders, Nick Young, Aaron Brooks \\
T3 & 444 & 0.96 & 10.6 & 13.0 & 3 & -0.86 & 1 & 3.0 & low-PER, short career, declining & Jonny Flynn, Anthony Jones, Litterial Green \\
T4 & 303 & 0.98 & 10.6 & 12.4 & 3 & +0.70 & 0 & 1.5 & low-PER, short career, rising & Qyntel Woods, Delaney Rudd, Larry Spriggs \\
\bottomrule\end{tabular}}
\end{table}
\begin{table}[h]\centering\small
\caption{Cross-tabulation of the eight fine types (rows) against the four macro-types (columns); \nestingPurity\% of players lie in the modal macro-type of their fine type (ARI between the two partitions \ariFourEight).}\label{tab:nesting}
\begin{tabular}{lrrrr}
\toprule
Macro-type & 1 & 2 & 3 & 4 \\
Fine type &  &  &  &  \\
\midrule
1 & 138 & 0 & 0 & 0 \\
2 & 101 & 30 & 35 & 0 \\
3 & 92 & 190 & 1 & 0 \\
4 & 38 & 45 & 5 & 53 \\
5 & 0 & 12 & 224 & 40 \\
6 & 0 & 283 & 50 & 6 \\
7 & 0 & 0 & 2 & 204 \\
8 & 0 & 0 & 127 & 0 \\
\bottomrule
\end{tabular}

\end{table}

\section{Internal indices in each method's own space}\label{app:internal}
Table~\ref{tab:internal} reports Silhouette, Calinski--Harabasz and Davies--Bouldin indices of each $k=8$ consensus partition in the space that the respective method clustered; they are not comparable across representations.
\begin{table}[h]\centering\small
\caption{Internal indices of the $k=8$ consensus partitions, computed in each representation's own (standardised) space.}\label{tab:internal}
\begin{tabular}{lccc}\toprule
Representation & Silhouette & Calinski--Harabasz & Davies--Bouldin \\ \midrule
Summary statistics, $k$-means & 0.19 & 395 & 1.40 \\
Resampled trajectory, $k$-means & 0.17 & 548 & 1.53 \\
PCA of resampled trajectory, $k$-means & 0.02 & 145 & 2.63 \\
Autoencoder $d=64$, 5 seeds (latent space) & 0.09 & 127 & 2.24 \\
\bottomrule\end{tabular}
\end{table}
